\section{Controller Design}

\subsection{Nominal Controller Design}

To achieve the desired tracking objective, define the auxiliary control input as

\begin{equation}
	v
	=
	\ddot q_d
	-
	K_v\dot e
	-
	K_pe,
	\label{eq:v}
\end{equation}

where $K_p,K_v\in\mathbb{R}^{n\times n}$ are symmetric positive-definite gain matrices.

Based on the nominal manipulator dynamics, the control law is designed as

\begin{equation}
	\tau
	=
	M_0(q)v
	+
	C_0(q,\dot q)\dot q
	+
	G_0(q)
	-
	\hat f(x),
	\label{eq:controller}
\end{equation}

where $\hat f(x)$ denotes the neural-network estimation of the lumped uncertainty.

\begin{remark}
	The introduction of the virtual control input $v$ decouples the controller into two functional modules. The outer-loop controller computes the virtual acceleration command $v$ to achieve the desired tracking performance, whereas the inner-loop controller realizes this command by generating the actual joint torque $\tau$ through dynamic inversion and RBF neural network compensation. This modular design separates trajectory tracking from dynamic compensation, thereby simplifying controller design and implementation.
\end{remark}

Substituting \eqref{eq:controller} into the nominal dynamics \eqref{eq:nominal} yields

\begin{equation}
	M_0(q)\ddot q
	=
	M_0(q)v
	+
	f(q,\dot q,\ddot q)
	-
	\hat f(x).
\end{equation}

Using the definition of $v$ in \eqref{eq:v}, one obtains

\begin{equation}
	M_0(q)
	\left(
	\ddot q
	-
	\ddot q_d
	+
	K_v\dot e
	+
	K_pe
	\right)
	=
	f
	-
	\hat f(x).
\end{equation}

Since

\begin{equation}
	\ddot e
	=
	\ddot q
	-
	\ddot q_d,
\end{equation}

the closed-loop tracking error dynamics can be written as

\begin{equation}
	\ddot e
	+
	K_v\dot e
	+
	K_pe
	=
	M_0^{-1}(q)
	\left(
	f
	-
	\hat f(x)
	\right).
	\label{eq:closedloop}
\end{equation}

%%%%%%%%%%%%%%%%%%%%%%%%%%%%%%%%%%%%%%%%%%%%%%%%%%%%%%%%%%%%%%

\subsection{State-Space Representation}

Define the state vector as

\begin{equation}
	X
	=
	\begin{bmatrix}
		e\\
		\dot e
	\end{bmatrix},
	\qquad
	\dot X
	=
	\begin{bmatrix}
		\dot e\\
		\ddot e
	\end{bmatrix}.
	\label{eq:state}
\end{equation}

Substituting \eqref{eq:closedloop} into \eqref{eq:state} gives

\begin{equation}
	\dot X
	=
	AX
	+
	B
	\left(
	f
	-
	\hat f(x)
	\right),
	\label{eq:ss1}
\end{equation}

where

\begin{equation}
	A
	=
	\begin{bmatrix}
		0&I\\
		-K_p&-K_v
	\end{bmatrix},
	\qquad
	B
	=
	\begin{bmatrix}
		0\\
		M_0^{-1}(q)
	\end{bmatrix}.
	\label{eq:AB}
\end{equation}

Equation \eqref{eq:ss1} reformulates the second-order closed-loop system as a first-order state-space representation, facilitating the subsequent stability analysis and adaptive controller design.

%%%%%%%%%%%%%%%%%%%%%%%%%%%%%%%%%%%%%%%%%%%%%%%%%%%%%%%%%%%%%%

\subsection{RBF Neural Network Approximation}

To compensate for the unknown lumped uncertainty $f(x)$, an RBF neural network is employed.

According to the universal approximation theorem, there exists an ideal constant weight matrix $W^*$ such that

\begin{equation}
	f(x)
	=
	W^{*T}h(x)
	+
	\varepsilon,
	\label{eq:rbf}
\end{equation}

where $h(x)\in\mathbb{R}^{m}$ denotes the Gaussian basis function vector, $W^*\in\mathbb{R}^{m\times n}$ is the ideal weight matrix, and $\varepsilon\in\mathbb{R}^{n}$ represents the bounded approximation error satisfying

\begin{equation}
	\|\varepsilon\|
	\le
	\bar\varepsilon.
\end{equation}

The neural-network output is given by

\begin{equation}
	\hat f(x)
	=
	\hat W^Th(x),
	\label{eq:fhat}
\end{equation}

where $\hat W$ denotes the online estimate of the ideal weight matrix.

Define the weight estimation error as

\begin{equation}
	\tilde W
	=
	W^*
	-
	\hat W.
	\label{eq:weight}
\end{equation}

Substituting \eqref{eq:rbf} and \eqref{eq:fhat} into \eqref{eq:ss1} yields

\begin{equation}
	\dot X
	=
	AX
	+
	B
	\left(
	\tilde W^Th(x)
	+
	\varepsilon
	\right),
	\label{eq:ss}
\end{equation}

which will be employed in the subsequent Lyapunov stability analysis.